% --- Archon physics formalization source begin ---
% archon:physics
% archon:covers PhyXMiniProblems/problem_phyx_mini_0551.lean
% archon:source-report reports/phyx_mini/problem_phyx_mini_0551.source.json

\chapter{Physics problem phyx\_mini\_0551}
\label{ch:PhyXMiniProblems_problem_phyx_mini_0551}

\paragraph{Problem source.}
\#\# Physical scenario

An atom of uranium (\textbackslash{}( m = 3.9529 \textbackslash{}times 10\textasciicircum{}\{-25\} \textbackslash{}) kg) at rest decays spontaneously into an atom of helium (\textbackslash{}( m = 6.6465 \textbackslash{}times 10\textasciicircum{}\{-27\} \textbackslash{}) kg) and an atom of thorium (\textbackslash{}( m = 3.8864 \textbackslash{}times 10\textasciicircum{}\{-25\} \textbackslash{}) kg). The helium atom is observed to move in the positive \textbackslash{}( x \textbackslash{}) direction with a velocity of \textbackslash{}( 1.423 \textbackslash{}times 10\textasciicircum{}7 \textbackslash{}) m/s as shown in figure.

\#\# Question

Find the total kinetic energy of the two atoms after the decay.

\#\# Answer choices

- A: A: \textbackslash{}( 6.09 \textbackslash{}times 10\textasciicircum{}\{-13\} \textbackslash{}text\{ J\} \textbackslash{})

- B: B: \textbackslash{}( 5.992 \textbackslash{}times 10\textasciicircum{}\{-13\} \textbackslash{}text\{ J\} \textbackslash{})

- C: C: \textbackslash{}( 7.231 \textbackslash{}times 10\textasciicircum{}\{-13\} \textbackslash{}text\{ J\} \textbackslash{})

- D: D: \textbackslash{}( 6.844 \textbackslash{}times 10\textasciicircum{}\{-13\} \textbackslash{}text\{ J\} \textbackslash{})

\#\# Figure caption (auxiliary; use the image as primary evidence)

The image is divided into two panels, both containing diagrams with labeled axes and particles.

**Top Panel:**
- A single sphere is positioned at the origin of an x-y coordinate system.
- The sphere is labeled "U," representing Uranium.

**Bottom Panel:**
- Two spheres are depicted on an x-axis. 
- The left sphere, labeled "Th," represents Thorium and has an arrow pointing left, labeled "ν'\_Th."
- The right sphere, labeled "He," represents Helium and has an arrow pointing right, labeled "ν'\_He."
- Both spheres are aligned horizontally on the x-axis.

The diagrams illustrate a nuclear reaction involving Uranium splitting into Thorium and Helium, with corresponding velocities.

\#\# Dataset metadata

Nuclear Physics; Multi-Formula Reasoning, Physical Model Grounding Reasoning

\paragraph{Recorded answer/context.}
D: D: \textbackslash{}( 6.844 \textbackslash{}times 10\textasciicircum{}\{-13\} \textbackslash{}text\{ J\} \textbackslash{})

\paragraph{Figure/image path.}
/root/proposal\_for\_physic/science-mango/phyx\_mini\_run/phyx\_data/test\_image/551.png

\paragraph{Formalization target.}
create a compiling Lean file with sorry bodies at `PhyXMiniProblems/problem\_phyx\_mini\_0551.lean`. The Lean declarations must preserve the physical quantities, dimensions or dimensional roles, figure labels, governing-law hypotheses, and final relation expressed by this problem.
Use Mathlib/Physlib names found through LeanExplore where available. If a physics API is missing, introduce faithful local abstractions rather than scalar placeholder aliases.

\begin{theorem}[Physics formalization target]
\label{thm:physics:phyx_mini_0551:target}
\lean{PhyXMiniProblems.ProblemPhyXMini0551.problem_phyx_mini_0551}
\uses{def:physics:phyx-mini-0551:phyxminiproblems-problemphyxmini0551-energyinjoules, def:physics:phyx-mini-0551:phyxminiproblems-problemphyxmini0551-uraniumdecaysetup, def:physics:phyx-mini-0551:phyxminiproblems-problemphyxmini0551-matchesuraniumdecayscenario, def:physics:phyx-mini-0551:phyxminiproblems-problemphyxmini0551-matchesuraniumdecayreadouts, def:physics:phyx-mini-0551:phyxminiproblems-problemphyxmini0551-hasphysicaluraniumdecayparameters, def:physics:phyx-mini-0551:phyxminiproblems-problemphyxmini0551-matchessupplieduraniumdecayfigure, def:physics:phyx-mini-0551:phyxminiproblems-problemphyxmini0551-satisfiesfigurevelocitydirectionconvention, def:physics:phyx-mini-0551:phyxminiproblems-problemphyxmini0551-satisfiesdecaymomentumconservation, def:physics:phyx-mini-0551:phyxminiproblems-problemphyxmini0551-satisfiesnonrelativistickineticenergylaws, def:physics:phyx-mini-0551:phyxminiproblems-problemphyxmini0551-isuniqueclosestdisplayedenergy}
The assigned autoformalize agent should translate this physics problem into Lean declarations in the covered file, with theorem and lemma proof bodies written as `by sorry`.
\end{theorem}
\begin{proof}
This is an autoformalization task, not a proof task. Produce statements that can later be proved by the physics prover without weakening the physical model.
\end{proof}

% --- Archon declaration topology begin ---
\section{Declaration topology}
\begin{definition}[Mass Quantity]
  \label{def:physics:phyx-mini-0551:phyxminiproblems-problemphyxmini0551-massquantity}
  \lean{PhyXMiniProblems.ProblemPhyXMini0551.MassQuantity}
  A nonnegative physical mass, independent of the unit used to read it
\end{definition}
\begin{proof}
  This is the declared model definition of Mass Quantity.
\end{proof}
\begin{definition}[Signed Axial Velocity Quantity]
  \label{def:physics:phyx-mini-0551:phyxminiproblems-problemphyxmini0551-signedaxialvelocityquantity}
  \lean{PhyXMiniProblems.ProblemPhyXMini0551.SignedAxialVelocityQuantity}
  A signed one-dimensional physical velocity along the decay axis
\end{definition}
\begin{proof}
  This is the declared model definition of Signed Axial Velocity Quantity.
\end{proof}
\begin{definition}[Energy Quantity]
  \label{def:physics:phyx-mini-0551:phyxminiproblems-problemphyxmini0551-energyquantity}
  \lean{PhyXMiniProblems.ProblemPhyXMini0551.EnergyQuantity}
  A signed physical energy; kinetic energies are constrained to be nonnegative
\end{definition}
\begin{proof}
  This is the declared model definition of Energy Quantity.
\end{proof}
\begin{definition}[mass Readout]
  \label{def:physics:phyx-mini-0551:phyxminiproblems-problemphyxmini0551-massreadout}
  \lean{PhyXMiniProblems.ProblemPhyXMini0551.massReadout}
  \uses{def:physics:phyx-mini-0551:phyxminiproblems-problemphyxmini0551-massquantity}
  Read a physical mass in a selected mass unit
\end{definition}
\begin{proof}
  This is the declared model definition of mass Readout.
\end{proof}
\begin{definition}[velocity Readout]
  \label{def:physics:phyx-mini-0551:phyxminiproblems-problemphyxmini0551-velocityreadout}
  \lean{PhyXMiniProblems.ProblemPhyXMini0551.velocityReadout}
  \uses{def:physics:phyx-mini-0551:phyxminiproblems-problemphyxmini0551-signedaxialvelocityquantity}
  Read a signed axial velocity in coherent selected length and time units
\end{definition}
\begin{proof}
  This is the declared model definition of velocity Readout.
\end{proof}
\begin{definition}[energy Readout]
  \label{def:physics:phyx-mini-0551:phyxminiproblems-problemphyxmini0551-energyreadout}
  \lean{PhyXMiniProblems.ProblemPhyXMini0551.energyReadout}
  \uses{def:physics:phyx-mini-0551:phyxminiproblems-problemphyxmini0551-energyquantity}
  Read an energy in the coherent unit determined by the selected base units
\end{definition}
\begin{proof}
  This is the declared model definition of energy Readout.
\end{proof}
\begin{definition}[mass In Kilograms]
  \label{def:physics:phyx-mini-0551:phyxminiproblems-problemphyxmini0551-massinkilograms}
  \lean{PhyXMiniProblems.ProblemPhyXMini0551.massInKilograms}
  \uses{def:physics:phyx-mini-0551:phyxminiproblems-problemphyxmini0551-massquantity, def:physics:phyx-mini-0551:phyxminiproblems-problemphyxmini0551-massreadout}
  Kilogram readout of a physical mass
\end{definition}
\begin{proof}
  This is the declared model definition of mass In Kilograms.
\end{proof}
\begin{definition}[velocity In Meters Per Second]
  \label{def:physics:phyx-mini-0551:phyxminiproblems-problemphyxmini0551-velocityinmeterspersecond}
  \lean{PhyXMiniProblems.ProblemPhyXMini0551.velocityInMetersPerSecond}
  \uses{def:physics:phyx-mini-0551:phyxminiproblems-problemphyxmini0551-signedaxialvelocityquantity, def:physics:phyx-mini-0551:phyxminiproblems-problemphyxmini0551-velocityreadout}
  Metre-per-second readout of a signed axial velocity
\end{definition}
\begin{proof}
  This is the declared model definition of velocity In Meters Per Second.
\end{proof}
\begin{definition}[energy In Joules]
  \label{def:physics:phyx-mini-0551:phyxminiproblems-problemphyxmini0551-energyinjoules}
  \lean{PhyXMiniProblems.ProblemPhyXMini0551.energyInJoules}
  \uses{def:physics:phyx-mini-0551:phyxminiproblems-problemphyxmini0551-energyquantity, def:physics:phyx-mini-0551:phyxminiproblems-problemphyxmini0551-energyreadout}
  Joule readout of a physical energy
\end{definition}
\begin{proof}
  This is the declared model definition of energy In Joules.
\end{proof}
\begin{definition}[Atom Species]
  \label{def:physics:phyx-mini-0551:phyxminiproblems-problemphyxmini0551-atomspecies}
  \lean{PhyXMiniProblems.ProblemPhyXMini0551.AtomSpecies}
  The three atom labels in the physical process and its figure
\end{definition}
\begin{proof}
  This is the declared model definition of Atom Species.
\end{proof}
\begin{definition}[Decay Product]
  \label{def:physics:phyx-mini-0551:phyxminiproblems-problemphyxmini0551-decayproduct}
  \lean{PhyXMiniProblems.ProblemPhyXMini0551.DecayProduct}
  The two products of the spontaneous decay
\end{definition}
\begin{proof}
  This is the declared model definition of Decay Product.
\end{proof}
\begin{definition}[atom Species]
  \label{def:physics:phyx-mini-0551:phyxminiproblems-problemphyxmini0551-decayproduct-atomspecies}
  \lean{PhyXMiniProblems.ProblemPhyXMini0551.DecayProduct.atomSpecies}
  \uses{def:physics:phyx-mini-0551:phyxminiproblems-problemphyxmini0551-atomspecies, def:physics:phyx-mini-0551:phyxminiproblems-problemphyxmini0551-decayproduct}
  Regard a decay product as its corresponding atom species
\end{definition}
\begin{proof}
  This is the declared model definition of atom Species.
\end{proof}
\begin{definition}[Coordinate Axis]
  \label{def:physics:phyx-mini-0551:phyxminiproblems-problemphyxmini0551-coordinateaxis}
  \lean{PhyXMiniProblems.ProblemPhyXMini0551.CoordinateAxis}
  The coordinate axes drawn in both panels of the source image
\end{definition}
\begin{proof}
  This is the declared model definition of Coordinate Axis.
\end{proof}
\begin{definition}[Figure Panel]
  \label{def:physics:phyx-mini-0551:phyxminiproblems-problemphyxmini0551-figurepanel}
  \lean{PhyXMiniProblems.ProblemPhyXMini0551.FigurePanel}
  The states shown above and below the horizontal divider in the image
\end{definition}
\begin{proof}
  This is the declared model definition of Figure Panel.
\end{proof}
\begin{definition}[Horizontal Location]
  \label{def:physics:phyx-mini-0551:phyxminiproblems-problemphyxmini0551-horizontallocation}
  \lean{PhyXMiniProblems.ProblemPhyXMini0551.HorizontalLocation}
  Qualitative horizontal locations relative to the drawn coordinate origin
\end{definition}
\begin{proof}
  This is the declared model definition of Horizontal Location.
\end{proof}
\begin{definition}[Horizontal Direction]
  \label{def:physics:phyx-mini-0551:phyxminiproblems-problemphyxmini0551-horizontaldirection}
  \lean{PhyXMiniProblems.ProblemPhyXMini0551.HorizontalDirection}
  Directions of the two post-decay velocity arrows
\end{definition}
\begin{proof}
  This is the declared model definition of Horizontal Direction.
\end{proof}
\begin{definition}[Decay Mode]
  \label{def:physics:phyx-mini-0551:phyxminiproblems-problemphyxmini0551-decaymode}
  \lean{PhyXMiniProblems.ProblemPhyXMini0551.DecayMode}
  The physical decay classification stated in the scenario
\end{definition}
\begin{proof}
  This is the declared model definition of Decay Mode.
\end{proof}
\begin{definition}[Uranium Decay Figure]
  \label{def:physics:phyx-mini-0551:phyxminiproblems-problemphyxmini0551-uraniumdecayfigure}
  \lean{PhyXMiniProblems.ProblemPhyXMini0551.UraniumDecayFigure}
  \uses{def:physics:phyx-mini-0551:phyxminiproblems-problemphyxmini0551-atomspecies, def:physics:phyx-mini-0551:phyxminiproblems-problemphyxmini0551-decayproduct, def:physics:phyx-mini-0551:phyxminiproblems-problemphyxmini0551-coordinateaxis, def:physics:phyx-mini-0551:phyxminiproblems-problemphyxmini0551-figurepanel, def:physics:phyx-mini-0551:phyxminiproblems-problemphyxmini0551-horizontallocation, def:physics:phyx-mini-0551:phyxminiproblems-problemphyxmini0551-horizontaldirection}
  Qualitative and raster evidence extracted from the supplied primary image. There is no quantitative energy or unobserved thorium speed in this structure
\end{definition}
\begin{proof}
  This is the declared model definition of Uranium Decay Figure.
\end{proof}
\begin{definition}[Uranium Decay Setup]
  \label{def:physics:phyx-mini-0551:phyxminiproblems-problemphyxmini0551-uraniumdecaysetup}
  \lean{PhyXMiniProblems.ProblemPhyXMini0551.UraniumDecaySetup}
  \uses{def:physics:phyx-mini-0551:phyxminiproblems-problemphyxmini0551-massquantity, def:physics:phyx-mini-0551:phyxminiproblems-problemphyxmini0551-signedaxialvelocityquantity, def:physics:phyx-mini-0551:phyxminiproblems-problemphyxmini0551-energyquantity, def:physics:phyx-mini-0551:phyxminiproblems-problemphyxmini0551-decayproduct, def:physics:phyx-mini-0551:phyxminiproblems-problemphyxmini0551-coordinateaxis, def:physics:phyx-mini-0551:phyxminiproblems-problemphyxmini0551-decaymode, def:physics:phyx-mini-0551:phyxminiproblems-problemphyxmini0551-uraniumdecayfigure}
  All independent physical quantities in the modeled event. In particular, the thorium velocity, both daughter kinetic energies, and their total are unconstrained fields until the governing laws are supplied; none is defined from an answer choice
\end{definition}
\begin{proof}
  This is the declared model definition of Uranium Decay Setup.
\end{proof}
\begin{definition}[Matches Uranium Decay Scenario]
  \label{def:physics:phyx-mini-0551:phyxminiproblems-problemphyxmini0551-matchesuraniumdecayscenario}
  \lean{PhyXMiniProblems.ProblemPhyXMini0551.MatchesUraniumDecayScenario}
  \uses{def:physics:phyx-mini-0551:phyxminiproblems-problemphyxmini0551-uraniumdecaysetup}
  The qualitative spontaneous, one-dimensional decay scenario
\end{definition}
\begin{proof}
  This is the declared model definition of Matches Uranium Decay Scenario.
\end{proof}
\begin{definition}[Matches Uranium Decay Readouts]
  \label{def:physics:phyx-mini-0551:phyxminiproblems-problemphyxmini0551-matchesuraniumdecayreadouts}
  \lean{PhyXMiniProblems.ProblemPhyXMini0551.MatchesUraniumDecayReadouts}
  \uses{def:physics:phyx-mini-0551:phyxminiproblems-problemphyxmini0551-massinkilograms, def:physics:phyx-mini-0551:phyxminiproblems-problemphyxmini0551-velocityinmeterspersecond, def:physics:phyx-mini-0551:phyxminiproblems-problemphyxmini0551-uraniumdecaysetup}
  The three stated masses, the initially stationary uranium atom, and the observed signed helium velocity. No thorium velocity or energy is assumed
\end{definition}
\begin{proof}
  This is the declared model definition of Matches Uranium Decay Readouts.
\end{proof}
\begin{definition}[Has Physical Uranium Decay Parameters]
  \label{def:physics:phyx-mini-0551:phyxminiproblems-problemphyxmini0551-hasphysicaluraniumdecayparameters}
  \lean{PhyXMiniProblems.ProblemPhyXMini0551.HasPhysicalUraniumDecayParameters}
  \uses{def:physics:phyx-mini-0551:phyxminiproblems-problemphyxmini0551-massinkilograms, def:physics:phyx-mini-0551:phyxminiproblems-problemphyxmini0551-energyinjoules, def:physics:phyx-mini-0551:phyxminiproblems-problemphyxmini0551-uraniumdecaysetup}
  Positivity and nonnegativity conditions selecting the physical branch
\end{definition}
\begin{proof}
  This is the declared model definition of Has Physical Uranium Decay Parameters.
\end{proof}
\begin{definition}[Matches Supplied Uranium Decay Figure]
  \label{def:physics:phyx-mini-0551:phyxminiproblems-problemphyxmini0551-matchessupplieduraniumdecayfigure}
  \lean{PhyXMiniProblems.ProblemPhyXMini0551.MatchesSuppliedUraniumDecayFigure}
  \uses{def:physics:phyx-mini-0551:phyxminiproblems-problemphyxmini0551-uraniumdecayfigure}
  Visible contents of the 430 × 400 source raster. The uranium atom occupies the origin in the upper panel. In the lower panel, thorium lies to the left and helium to the right, with arrows labeled v'\_Th and v'\_He pointing in the corresponding directions
\end{definition}
\begin{proof}
  This is the declared model definition of Matches Supplied Uranium Decay Figure.
\end{proof}
\begin{definition}[Satisfies Figure Velocity Direction Convention]
  \label{def:physics:phyx-mini-0551:phyxminiproblems-problemphyxmini0551-satisfiesfigurevelocitydirectionconvention}
  \lean{PhyXMiniProblems.ProblemPhyXMini0551.SatisfiesFigureVelocityDirectionConvention}
  \uses{def:physics:phyx-mini-0551:phyxminiproblems-problemphyxmini0551-velocityinmeterspersecond, def:physics:phyx-mini-0551:phyxminiproblems-problemphyxmini0551-uraniumdecaysetup}
  The conventional interpretation of each displayed velocity arrow. It links qualitative figure direction to the sign of a physical velocity but supplies neither the thorium speed magnitude nor any energy
\end{definition}
\begin{proof}
  This is the declared model definition of Satisfies Figure Velocity Direction Convention.
\end{proof}
\begin{definition}[Satisfies Decay Momentum Conservation]
  \label{def:physics:phyx-mini-0551:phyxminiproblems-problemphyxmini0551-satisfiesdecaymomentumconservation}
  \lean{PhyXMiniProblems.ProblemPhyXMini0551.SatisfiesDecayMomentumConservation}
  \uses{def:physics:phyx-mini-0551:phyxminiproblems-problemphyxmini0551-massreadout, def:physics:phyx-mini-0551:phyxminiproblems-problemphyxmini0551-velocityreadout, def:physics:phyx-mini-0551:phyxminiproblems-problemphyxmini0551-uraniumdecaysetup}
  One-dimensional linear momentum conservation for the isolated two-body decay, stated in every coherent choice of mass, length, and time units. The parent term is retained explicitly so the initially-at-rest datum remains separate
\end{definition}
\begin{proof}
  This is the declared model definition of Satisfies Decay Momentum Conservation.
\end{proof}
\begin{definition}[Satisfies Nonrelativistic Kinetic Energy Laws]
  \label{def:physics:phyx-mini-0551:phyxminiproblems-problemphyxmini0551-satisfiesnonrelativistickineticenergylaws}
  \lean{PhyXMiniProblems.ProblemPhyXMini0551.SatisfiesNonrelativisticKineticEnergyLaws}
  \uses{def:physics:phyx-mini-0551:phyxminiproblems-problemphyxmini0551-massreadout, def:physics:phyx-mini-0551:phyxminiproblems-problemphyxmini0551-velocityreadout, def:physics:phyx-mini-0551:phyxminiproblems-problemphyxmini0551-energyreadout, def:physics:phyx-mini-0551:phyxminiproblems-problemphyxmini0551-decayproduct, def:physics:phyx-mini-0551:phyxminiproblems-problemphyxmini0551-uraniumdecaysetup}
  The nonrelativistic translational kinetic-energy law for each daughter and the additivity of their total kinetic energy. These are general governing laws; they contain no measured or displayed numerical energy
\end{definition}
\begin{proof}
  This is the declared model definition of Satisfies Nonrelativistic Kinetic Energy Laws.
\end{proof}
\begin{lemma}[thorium recoil velocity exact]
  \label{lem:physics:phyx-mini-0551:phyxminiproblems-problemphyxmini0551-thorium-recoil-velocity-exact}
  \lean{PhyXMiniProblems.ProblemPhyXMini0551.thorium_recoil_velocity_exact}
  \uses{def:physics:phyx-mini-0551:phyxminiproblems-problemphyxmini0551-massinkilograms, def:physics:phyx-mini-0551:phyxminiproblems-problemphyxmini0551-velocityinmeterspersecond, def:physics:phyx-mini-0551:phyxminiproblems-problemphyxmini0551-uraniumdecaysetup, def:physics:phyx-mini-0551:phyxminiproblems-problemphyxmini0551-matchesuraniumdecayreadouts, def:physics:phyx-mini-0551:phyxminiproblems-problemphyxmini0551-hasphysicaluraniumdecayparameters, def:physics:phyx-mini-0551:phyxminiproblems-problemphyxmini0551-satisfiesdecaymomentumconservation}
  Momentum conservation determines the signed thorium recoil velocity. This is a derived conclusion and is not a premise of the main theorem
\end{lemma}
\begin{proof}
  The conclusion follows from the governing relations and figure readouts named in the statement of thorium recoil velocity exact.
\end{proof}
\begin{lemma}[total kinetic energy exact]
  \label{lem:physics:phyx-mini-0551:phyxminiproblems-problemphyxmini0551-total-kinetic-energy-exact}
  \lean{PhyXMiniProblems.ProblemPhyXMini0551.total_kinetic_energy_exact}
  \uses{def:physics:phyx-mini-0551:phyxminiproblems-problemphyxmini0551-massinkilograms, def:physics:phyx-mini-0551:phyxminiproblems-problemphyxmini0551-velocityinmeterspersecond, def:physics:phyx-mini-0551:phyxminiproblems-problemphyxmini0551-energyinjoules, def:physics:phyx-mini-0551:phyxminiproblems-problemphyxmini0551-uraniumdecaysetup, def:physics:phyx-mini-0551:phyxminiproblems-problemphyxmini0551-matchesuraniumdecayreadouts, def:physics:phyx-mini-0551:phyxminiproblems-problemphyxmini0551-hasphysicaluraniumdecayparameters, def:physics:phyx-mini-0551:phyxminiproblems-problemphyxmini0551-satisfiesdecaymomentumconservation, def:physics:phyx-mini-0551:phyxminiproblems-problemphyxmini0551-satisfiesnonrelativistickineticenergylaws}
  Substituting the recoil relation into K = m v² / 2 gives an exact SI formula for the total daughter kinetic energy before numerical rounding
\end{lemma}
\begin{proof}
  The conclusion follows from the governing relations and figure readouts named in the statement of total kinetic energy exact.
\end{proof}
\begin{definition}[Answer Choice]
  \label{def:physics:phyx-mini-0551:phyxminiproblems-problemphyxmini0551-answerchoice}
  \lean{PhyXMiniProblems.ProblemPhyXMini0551.AnswerChoice}
  Labels printed beside the four answer choices
\end{definition}
\begin{proof}
  This is the declared model definition of Answer Choice.
\end{proof}
\begin{definition}[energy In Joules]
  \label{def:physics:phyx-mini-0551:phyxminiproblems-problemphyxmini0551-answerchoice-energyinjoules}
  \lean{PhyXMiniProblems.ProblemPhyXMini0551.AnswerChoice.energyInJoules}
  \uses{def:physics:phyx-mini-0551:phyxminiproblems-problemphyxmini0551-energyinjoules, def:physics:phyx-mini-0551:phyxminiproblems-problemphyxmini0551-answerchoice}
  Total kinetic energy printed beside each answer label, in joules
\end{definition}
\begin{proof}
  This is the declared model definition of energy In Joules.
\end{proof}
\begin{definition}[recorded Dataset Answer]
  \label{def:physics:phyx-mini-0551:phyxminiproblems-problemphyxmini0551-recordeddatasetanswer}
  \lean{PhyXMiniProblems.ProblemPhyXMini0551.recordedDatasetAnswer}
  \uses{def:physics:phyx-mini-0551:phyxminiproblems-problemphyxmini0551-answerchoice}
  The answer label recorded in the source dataset metadata
\end{definition}
\begin{proof}
  This is the declared model definition of recorded Dataset Answer.
\end{proof}
\begin{definition}[Is Unique Closest Displayed Energy]
  \label{def:physics:phyx-mini-0551:phyxminiproblems-problemphyxmini0551-isuniqueclosestdisplayedenergy}
  \lean{PhyXMiniProblems.ProblemPhyXMini0551.IsUniqueClosestDisplayedEnergy}
  \uses{def:physics:phyx-mini-0551:phyxminiproblems-problemphyxmini0551-energyinjoules, def:physics:phyx-mini-0551:phyxminiproblems-problemphyxmini0551-uraniumdecaysetup, def:physics:phyx-mini-0551:phyxminiproblems-problemphyxmini0551-answerchoice}
  A displayed choice is uniquely closest to the physical total-energy readout. This is a target-side comparison, not an assumed law or data calibration
\end{definition}
\begin{proof}
  This is the declared model definition of Is Unique Closest Displayed Energy.
\end{proof}
% --- Archon declaration topology end ---
% --- Archon physics formalization source end ---
